% Part: first-order-logic
% Chapter: syntax-and-semantics
% Section: Semantic Notions

\documentclass[../../../include/open-logic-section]{subfiles}

\begin{document}

\olfileid{fol}{syn}{sem}
\olsection{অর্থগত ধারণা}

\begin{explain}
প্রথম-ক্রমের ভাষার !!{structure}s-এর সংজ্ঞা দেওয়া থাকলে বাক্যের কয়েকটি
মৌলিক অর্থগত ধর্ম এবং বাক্যগুলির মধ্যেকার সম্বন্ধ সংজ্ঞায়িত করতে পারি।
এর মধ্যে সবচেয়ে সরল হলো বাক্যের \emph{বৈধতা}। কোনো বাক্য প্রত্যেক
!!{structure}-এ পরিতৃপ্ত হলে সেটি বৈধ। বৈধ বাক্য তার অ-যৌক্তিক
সংকেতগুলিকে যেভাবেই ব্যাখ্যা করা হোক না কেন পরিতৃপ্ত হয়। তাই বৈধ
বাক্যকে \emph{যৌক্তিক সত্য}-ও বলা হয়—সেগুলি যেকোনো !!{structure}-এ
সত্য, অর্থাৎ পরিতৃপ্ত; ফলে তাদের সত্যতা শুধু তাদের মধ্যে সংঘটিত যৌক্তিক
সংকেত এবং তাদের সংকেতবিন্যাসগত !!{structure}-এর উপর নির্ভর করে,
অ-যৌক্তিক সংকেত বা তার ব্যাখ্যার উপর নয়।
\end{explain}

\begin{defn}[বৈধতা]
কোনো বাক্য $!A$ \emph{বৈধ}, $\Entails !A$, যদি এবং কেবল যদি প্রত্যেক
!!{structure}~$\Struct M$-এর জন্য $\Sat{M}{!A}$।
\end{defn}

\begin{defn}[ফলশ্রুতি]
বাক্যসমষ্টি~$\Gamma$ থেকে বাক্য~$!A$ \emph{অনুসৃত হয়},
$\Gamma \Entails !A$, যদি এবং কেবল যদি $\Sat{M}{\Gamma}$ এমন প্রত্যেক
!!{structure}~$\Struct M$-এর জন্য $\Sat{M}{!A}$।
\end{defn}


\begin{defn}[পরিতৃপ্তিযোগ্যতা]
বাক্যসমষ্টি~$\Gamma$ \emph{পরিতৃপ্তিযোগ্য}, যদি কোনো
!!{structure}~$\Struct M$-এর জন্য $\Sat{M}{\Gamma}$। $\Gamma$
পরিতৃপ্তিযোগ্য না হলে তাকে \emph{অপরিতৃপ্তিযোগ্য} বলে।
\end{defn}

\begin{prop}
কোনো বাক্য $!A$ তখনই বৈধ, যখন প্রত্যেক বাক্যসমষ্টি~$\Gamma$-র জন্য
$\Gamma \Entails !A$।
\end{prop}

\begin{proof}
সম্মুখ দিকের জন্য ধরি $!A$ বৈধ এবং $\Gamma$ একটি বাক্যসমষ্টি। এমন
!!a{structure}~$\Struct M$ নিই যাতে $\Sat{M}{\Gamma}$। $!A$ বৈধ বলে
$\Sat{M}{!A}$; অতএব $\Gamma \Entails !A$।

বিপরীত দিকের বিপরীত-প্রতিপাদ্যের জন্য ধরি $!A$ অবৈধ; তাই এমন
!!a{structure}~$\Struct M$ আছে যাতে $\Sat/{M}{!A}$। যখন
$\Gamma = \{ \ltrue \}$, তখন $\ltrue$ বৈধ বলে $\Sat{M}{\Gamma}$।
অতএব এমন !!a{structure}~$\Struct M$ আছে যাতে $\Sat{M}{\Gamma}$,
কিন্তু $\Sat/{M}{!A}$; তাই $\Gamma$ থেকে $!A$ অনুসৃত হয় না।
\end{proof}

\begin{prop}
\ollabel{prop:entails-unsat}
$\Gamma \Entails !A$ তখনই, যখন $\Gamma \cup \{\lnot !A\}$
অপরিতৃপ্তিযোগ্য।
\end{prop}

\begin{proof}
সম্মুখ দিকের জন্য ধরি $\Gamma \Entails !A$ এবং বিপরীত অনুমান করি যে,
এমন !!a{structure}~$\Struct M$ আছে যাতে $\Sat{M}{\Gamma
  \cup \{ \lnot !A \}}$। যেহেতু $\Sat{M}{\Gamma}$ এবং
$\Gamma \Entails !A$, তাই $\Sat{M}{!A}$। আবার
$\Sat{M}{\Gamma\cup \{ \lnot !A \}}$ বলে $\Sat{M}{\lnot !A}$;
সুতরাং $\Sat{M}{!A}$ ও $\Sat/{M}{!A}$ উভয়ই পাই—এটি স্ববিরোধ। তাই
এমন কোনো !!{structure}~$\Struct M$ থাকতে পারে না; অতএব
$\Gamma \cup \{ \lnot !A \}$ অপরিতৃপ্তিযোগ্য।

বিপরীত দিকের জন্য ধরি $\Gamma \cup \{ \lnot !A \}$
অপরিতৃপ্তিযোগ্য। তাই প্রত্যেক !!{structure}~$\Struct M$-এর জন্য হয়
$\Sat/{M}{\Gamma}$, নয়তো $\Sat{M}{!A}$। অতএব $\Sat{M}{\Gamma}$ এমন
প্রত্যেক !!{structure}~$\Struct M$-এর জন্য $\Sat{M}{!A}$; তাই
$\Gamma \Entails !A$।
\end{proof}

\begin{prob}
\begin{enumerate}
\item দেখান, $\Gamma \Entails \bot$ তখনই, যখন $\Gamma$ অপরিতৃপ্তিযোগ্য।
\item দেখান, $\Gamma \cup \{!A\} \Entails \bot$ তখনই, যখন
  $\Gamma \Entails \lnot !A$।
\item ধরি $c$,~$!A$ বা $\Gamma$-তে সংঘটিত হয় না। দেখান,
  $\Gamma \Entails \lforall[x][!A]$ তখনই, যখন $\Gamma \Entails
  \Subst{!A}{c}{x}$।
\end{enumerate}
\end{prob}

\begin{prop}
$\Gamma \subseteq \Gamma'$ এবং $\Gamma \Entails !A$ হলে
$\Gamma' \Entails !A$।
\end{prop}

\begin{proof}
ধরি $\Gamma \subseteq \Gamma'$ এবং $\Gamma \Entails !A$। এমন একটি
গঠন $\Struct M$ নিই যাতে $\Sat{M}{\Gamma'}$; তাহলে
$\Sat{M}{\Gamma}$, এবং $\Gamma \Entails !A$ বলে $\Sat{M}{!A}$। তাই
যখনই $\Sat{M}{\Gamma'}$, তখন $\Sat{M}{!A}$; অতএব
$\Gamma' \Entails !A$।
\end{proof}


\begin{thm}[অর্থগত নিঃসরণ উপপাদ্য]
\ollabel{thm:sem-deduction}
$\Gamma \cup \{!A\} \Entails !B$ তখনই, যখন
$\Gamma \Entails !A \lif !B$।
\end{thm}

\begin{proof}
সম্মুখ দিকের জন্য ধরি $\Gamma \cup \{ !A \} \Entails !B$ এবং এমন
!!a{structure}~$\Struct M$ নিই যাতে $\Sat{M}{\Gamma}$।
$\Sat{M}{!A}$ হলে $\Sat{M}{\Gamma \cup \{ !A \} }$; তাই
$\Gamma \cup \{ !A \}$ থেকে $!B$ অনুসৃত হয় বলে $\Sat{M}{!B}$।
অতএব $\Sat{M}{!A \lif !B}$; তাই $\Gamma \Entails !A \lif !B$।

বিপরীত দিকের জন্য ধরি $\Gamma \Entails !A \lif !B$ এবং এমন
!!a{structure}~$\Struct M$ নিই যাতে $\Sat{M}{\Gamma \cup \{ !A
  \}}$। তাহলে $\Sat{M}{\Gamma}$, তাই $\Sat{M}{!A \lif !B}$; এবং
$\Sat{M}{!A}$ বলে $\Sat{M}{!B}$। অতএব যখনই
$\Sat{M}{\Gamma \cup \{ !A \} }$, তখন $\Sat{M}{!B}$; তাই
$\Gamma \cup \{ !A \} \Entails !B$।
\end{proof}

\begin{prop}\ollabel{prop:quant-terms}
  ধরি $\Struct{M}$ একটি !!a{structure}, $!A(x)$ একটি মুক্ত চলরাশি~$x$-সহ
  !!a{formula} এবং $t$ একটি বদ্ধ পদ। তাহলে:
  \begin{tagenumerate}{prvEx,prvAll}
    \tagitem{prvEx}{$!A(t) \Entails \lexists[x][!A(x)]$}{}
    \tagitem{prvAll}{$\lforall[x][!A(x)] \Entails !A(t)$}{}
  \end{tagenumerate}
\end{prop}

\begin{proof}
  \begin{tagenumerate}{prvEx,prvAll}
    \tagitem{prvEx}{%
      \iftag{probEx}{অনুশীলনী।}{ধরি $\Sat{M}{!A(t)}$। এমন একটি
        চলরাশি-আরোপ $s$ নিই যাতে $s(x) = \Value{t}{M}$। $!A(t)$ একটি
        !!a{sentence} বলে $\Sat{M}{!A(t)}[s]$।
        \olref[ext]{prop:ext-formulas} থেকে $\Sat{M}{!A(x)}[s]$।
        \olref[ass]{prop:sat-quant} থেকে
        $\Sat{M}{\lexists[x][!A(x)]}$।}}{}
    \tagitem{prvAll}{%
      \iftag{probAll}{অনুশীলনী।}{ধরি
        $\Sat{M}{\lforall[x][!A(x)]}$। এমন একটি চলরাশি-আরোপ $s$ নিই
        যাতে $s(x) = \Value{t}{M}$। \olref[ass]{prop:sat-quant} থেকে
        $\Sat{M}{!A(x)}[s]$। \olref[ext]{prop:ext-formulas} থেকে
        $\Sat{M}{!A(t)}[s]$। $!A(t)$ একটি !!a{sentence} বলে
        \olref[ass]{prop:sentence-sat-true} থেকে $\Sat{M}{!A(t)}$।}}{}
  \end{tagenumerate}
\end{proof}

\begin{probtag}{probEx,probAll}
  \olref[fol][syn][sem]{prop:quant-terms}-এর প্রমাণ সম্পূর্ণ করুন।
\end{probtag}

\end{document}
